% 公卫统计 · Einfart/plain 学习总笔记兼容层
% 视觉设置集中于本文件和 pku-notes.sty；原章节、题号、标签保持不变。
\usepackage{multirow,calc,tikz,pgfplots,placeins,pdflscape,microtype}
\usetikzlibrary{arrows.meta,positioning,calc,patterns,decorations.pathreplacing,shapes.geometric}
\pgfplotsset{compat=1.18}
\colorlet{PKUInk}{black}
\colorlet{Ink}{black}
\colorlet{Muted}{PKUMuted}
\colorlet{Primary}{PKURed}
\colorlet{Accent}{PKURed}
\colorlet{Rule}{black!22}
\colorlet{Paper}{PKUTint}
\colorlet{CodeBackground}{PKUCodeBg}
\colorlet{CodeRule}{black!22}
\colorlet{NavyBlue}{PKURed}
\colorlet{DarkRed}{PKURed}
\colorlet{ForestGreen}{black!65}
\colorlet{BoxGray}{black!3}
\colorlet{DefBlue}{PKURed!4}
\colorlet{ThmGold}{PKURed!3}
\colorlet{ExGreen}{black!3}
\colorlet{RemPurple}{black!3}
% 代码语义配色独立于装饰色。
\colorlet{CodeKeyword}{PKUCodePurple}
\colorlet{CodeFunction}{PKUCodeBlue}
\colorlet{CodeString}{PKUCodeString}
\colorlet{CodeComment}{PKUCodeGreen}
\definecolor{CodeConstant}{HTML}{805600}
\providecommand{\bm}[1]{\symbf{#1}}
\setlength{\emergencystretch}{3em}
\clubpenalty=10000\widowpenalty=10000\displaywidowpenalty=10000
\allowdisplaybreaks[2]
\setcounter{MaxMatrixCols}{14}
\setlength{\tabcolsep}{5pt}
\renewcommand{\arraystretch}{1.18}
\setlength{\heavyrulewidth}{.7pt}
\setlength{\lightrulewidth}{.35pt}
\setlength{\cmidrulewidth}{.3pt}
\captionsetup{font={small,pkublack},labelfont={bf,sf,pkublack},
  textfont={pkublack},labelsep=quad,skip=6pt,hypcap=false}
% Einfart 原生为 article。用 titlesec 增加 chapter 层，保留书稿结构。
\makeatletter
\def\ttll@part{-1}
\makeatother
\titleclass{\chapter}{top}[\part]
\newcounter{chapter}
\renewcommand{\thechapter}{\arabic{chapter}}
\counterwithin{section}{chapter}
\counterwithin{equation}{chapter}
\counterwithin{figure}{chapter}
\counterwithin{table}{chapter}
\providecommand{\theHchapter}{\arabic{chapter}}
\renewcommand{\thesection}{\thechapter.\arabic{section}}
\renewcommand{\chaptermark}[1]{\markboth{第\thechapter 章\quad #1}{}}
\titleformat{\chapter}[display]
  {\raggedright\sffamily\bfseries\color{black}}
  {\large\bfseries\color{black}第\,\thechapter\,章}{7pt}
  {\fontsize{23}{31}\selectfont\bfseries\color{black}#1}
  [\vspace{8pt}{\color{PKURed}\rule{\linewidth}{.7pt}}]
\titleformat{name=\chapter,numberless}[display]
  {\raggedright\sffamily\bfseries\color{black}}{}{0pt}
  {\fontsize{23}{31}\selectfont\bfseries\color{black}#1}
  [\phantomsection\vspace{8pt}{\color{PKURed}\rule{\linewidth}{.7pt}}]
\titlespacing*{\chapter}{0pt}{0pt}{22pt}
\titleformat{\section}
  {\sffamily\bfseries\Large\color{black}}{\thesection}{.7em}{#1}
\titleformat{\subsection}
  {\sffamily\bfseries\large\color{black}}{\thesubsection}{.65em}{#1}
\titleformat{\subsubsection}
  {\sffamily\bfseries\normalsize\color{black}}{\thesubsubsection}{.6em}{#1}
\setcounter{secnumdepth}{2}\setcounter{tocdepth}{1}
\pretocmd{\section}{\ifinner\else\Needspace{6\baselineskip}\fi}{}{}
\pretocmd{\subsection}{\ifinner\else\Needspace{4\baselineskip}\fi}{}{}
\newcommand{\frontmatter}{\clearpage\pagenumbering{Roman}}
\newcommand{\mainmatter}{\clearpage\pagenumbering{arabic}\setcounter{chapter}{0}}
\titlecontents{chapter}[3.2em]
  {\addvspace{.55\baselineskip}\PKUTOCFont}
  {\contentslabel[第\thecontentslabel 章]{3.2em}}{\hspace*{-3.2em}}
  {\titlerule*[.7em]{.}\contentspage}
\titlecontents{section}[5.4em]
  {\addvspace{.12\baselineskip}\small\PKUTOCFont}
  {\contentslabel{3em}}{\hspace*{-3em}}
  {\titlerule*[.7em]{.}\contentspage}
\makeatletter
\providecommand{\@chapapp}{章}
\renewcommand{\@pnumwidth}{2.5em}
\renewcommand{\@tocrmarg}{3.5em}
\def\toclevel@chapter{0}
\newcommand{\ZBTableOfContents}{\clearpage\chapter*{目录}%
  \markboth{目录}{}\begingroup\hypersetup{linkcolor=black}%
  \setlength{\parskip}{0pt}\@starttoc{toc}\endgroup}
\makeatother
% 页眉页脚：低对比辅助信息，页码常规字重。
\newcommand{\BookShortTitle}{\PKUShortTitle}
\newcommand{\ZBPageFurniture}{\fancyhf{}%
  \fancyhead[L]{\footnotesize\normalfont\sffamily\color{Muted}\nouppercase{\leftmark}}%
  \fancyhead[R]{\footnotesize\normalfont\sffamily\color{Muted}\BookShortTitle}%
  \fancyfoot[C]{\small\normalfont\sffamily\color{Muted}\thepage}%
  \renewcommand{\headrulewidth}{.3pt}\renewcommand{\footrulewidth}{0pt}}
\fancypagestyle{pku}{\ZBPageFurniture}
\fancypagestyle{plain}{\ZBPageFurniture}
\hypersetup{linkcolor=PKURed,urlcolor=PKURed,citecolor=PKURed}
\pdfstringdefDisableCommands{\def\color#1{}\def\sffamily{}\def\bfseries{}%
  \def\enspace{ }\def\quad{ }\def\texorpdfstring#1#2{#2}}
% 保留原正文环境，标题均在字体选择后显式设为黑色。
\tcbset{zbbase/.style={enhanced jigsaw,breakable,sharp corners,frame hidden,
  boxrule=0pt,before skip=9pt plus 2pt,after skip=9pt plus 2pt,pad at break*=1.5mm}}
\newtcolorbox{zbderive}[1]{zbbase,colback=PKURed!3,
  borderline west={1.5pt}{0pt}{PKURed},left=11pt,right=9pt,top=7pt,bottom=7pt,
  before upper={\setlength{\parskip}{3pt}{\sffamily\bfseries\color{black}#1}\par\smallskip}}
\newtcolorbox{zbnote}[1][批注]{zbbase,colback=white,
  borderline west={.8pt}{0pt}{black!35},left=9pt,right=2pt,top=2pt,bottom=2pt,
  fontupper=\small\color{Muted},
  before upper={{\sffamily\bfseries\color{black}#1}\enspace}}
\newtcolorbox{zbcheck}[1][说明]{zbbase,colback=PKURed!2,
  borderline west={1.2pt}{0pt}{PKURed},left=10pt,right=8pt,top=5pt,bottom=5pt,
  fontupper=\small\color{black},
  before upper={{\sffamily\bfseries\color{black}#1}\enspace}}
\newtcolorbox{keyidea}[1][]{zbbase,colback=PKURed!2,
  borderline west={1.3pt}{0pt}{PKURed},left=11pt,right=9pt,top=6pt,bottom=6pt,
  before upper={\ifx\relax#1\relax\else{\sffamily\bfseries\color{black}#1}\par\smallskip\fi}}
\newcommand{\dpart}[1]{\par\noindent{\sffamily\bfseries\color{black}#1}\quad\ignorespaces}
\let\question\relax
\let\answer\relax
\newcommand{\question}[1]{\par\Needspace{3\baselineskip}\smallskip\noindent
  {\sffamily\bfseries\color{black}问}\quad #1\par\smallskip}
\newcommand{\answer}[1]{\par\smallskip\noindent
  {\sffamily\bfseries\color{black}答}\quad #1\par\smallskip}
\newlist{choices}{enumerate}{1}
\setlist[choices]{label=\Alph*.,leftmargin=2.2em,itemsep=4pt,topsep=4pt,parsep=0pt,labelsep=.6em}
\newcounter{zbquestion}
\newcommand{\exsetno}{0}
\renewcommand{\thezbquestion}{\exsetno.\arabic{zbquestion}}
\renewcommand{\theHzbquestion}{\exsetno.\arabic{zbquestion}}
\newcommand{\ZBQuestion}[2]{\par\addvspace{12pt}\ifinner\else\Needspace{5\baselineskip}\fi
  \refstepcounter{zbquestion}\label{q:#1}\noindent
  {\sffamily\bfseries\color{black}题\thezbquestion}\quad
  {\sffamily\bfseries\color{black}#2}\par\nobreak\smallskip}
% unicode-math 在文档开始时把 Question 定义为双问号数学符号。
% 在所有字体/符号初始化完成后恢复原笔记的题面命令。
\AddToHook{begindocument/end}{\let\Question\ZBQuestion}
\newcommand{\Answer}[2]{\par\addvspace{12pt}\ifinner\else\Needspace{5\baselineskip}\fi\noindent
  {\sffamily\bfseries\color{black}解析\enspace 题\ref*{q:#1}}\quad
  {\sffamily\bfseries\color{black}#2}%
  \hfill{\footnotesize\sffamily\hyperref[q:#1]{\color{Muted}返回题目}}%
  \par\nobreak\smallskip\label{a:#1}}
\newcommand{\worklines}[1]{\par\smallskip{\color{Rule}\vbox{\offinterlineskip
  \foreach \n in {1,...,#1}{\vskip7.5mm\hrule height.3pt width\linewidth}}}\par\vspace{6pt}}
\DefineVerbatimEnvironment{verbatim}{Verbatim}{breaklines,fontsize=\small,formatcom=\xeCJKVerbAddon}
\DefineVerbatimEnvironment{Highlighting}{Verbatim}{breaklines,commandchars=\\\{\},fontsize=\small,formatcom=\xeCJKVerbAddon}
\providecommand{\tightlist}{\setlength{\itemsep}{0pt}\setlength{\parskip}{0pt}}
\providecommand{\pandocbounded}[1]{#1}
\providecommand{\passthrough}[1]{#1}
\AtBeginEnvironment{longtable}{\small\setlength{\tabcolsep}{4pt}\renewcommand{\arraystretch}{1.22}}
\setlength{\LTpre}{6pt}\setlength{\LTpost}{6pt}
\newcommand{\ExerciseNumber}[1]{\renewcommand{\exsetno}{#1}}
\newcommand{\ExerciseChapter}[1]{\clearpage\setcounter{zbquestion}{0}%
  \setcounter{secnumdepth}{0}\thispagestyle{plain}\phantomsection
  \addcontentsline{toc}{section}{\protect\numberline{}习题与解析\enspace #1}%
  \addtocontents{toc}{\protect\setcounter{tocdepth}{0}}%
  \markboth{第\exsetno 章\quad 习题与解析}{}%
  {\raggedright\sffamily\bfseries\color{black}%
   {\large\bfseries\color{black}第\,\exsetno\,章\enspace 习题与解析}\par\vspace{7pt}%
   {\fontsize{23}{31}\selectfont\bfseries\color{black}#1}\par\vspace{8pt}%
   {\color{PKURed}\rule{\linewidth}{.7pt}}\par}\vspace{22pt}%
  \pdfbookmark[1]{习题与解析：#1}{ex:\exsetno}}
\newcommand{\EndExercises}{\addtocontents{toc}{\protect\setcounter{tocdepth}{1}}\setcounter{secnumdepth}{2}}
% 简约封面：课程名、英文名、居中校徽、学习网站与日期。
\newcommand{\MakeCover}[2]{%
  \begin{titlepage}\thispagestyle{empty}\hypersetup{pageanchor=false}%
  \begin{tikzpicture}[remember picture,overlay]
    \node[align=center,text width=.88\paperwidth]
      at ([yshift=60mm]current page.center)
      {{\sffamily\fontsize{32}{43}\selectfont\bfseries\color{black}#1}};
    \node[align=center,text width=.88\paperwidth]
      at ([yshift=42mm]current page.center)
      {{\sffamily\fontsize{14}{20}\selectfont\color{Muted}#2}};
    \node at (current page.center)
      {\includegraphics[width=44mm]{../common/pku-emblem.jpg}};
    \node[align=center,text width=.88\paperwidth]
      at ([yshift=39mm]current page.south)
      {{\small\sffamily\color{Muted}学习网站：\url{https://fxt-gw-pb.github.io/public-health-statistics-notes/}}};
    \node at ([yshift=24mm]current page.south)
      {{\small\color{Muted}2026年10月}};
  \end{tikzpicture}%
  \null\end{titlepage}\hypersetup{pageanchor=true}}
% ---- 2026-10-06 润色补充：讲解框、运行结果、章末复习 ----
\DefineVerbatimEnvironment{routput}{Verbatim}{breaklines,breakanywhere,
  fontsize=\fontsize{8.6}{11.2}\selectfont,formatcom=\color{black!80}\xeCJKVerbAddon,
  frame=leftline,framerule=.8pt,rulecolor=\color{black!28},framesep=7pt,xleftmargin=12pt}
\newtcolorbox{sikao}[1][说明]{zbbase,colback=black!2.5,
  borderline west={1.2pt}{0pt}{black!50},left=10pt,right=8pt,top=6pt,bottom=6pt,
  before upper={\setlength{\parskip}{4pt}{\sffamily\bfseries\color{black}#1}\par\smallskip}}
\newtcolorbox{fuxi}[1][本章要点]{zbbase,colback=PKURed!2,
  borderline west={1.5pt}{0pt}{PKURed},left=11pt,right=9pt,top=7pt,bottom=7pt,
  before upper={\setlength{\parskip}{3pt}{\sffamily\bfseries\color{black}#1}\par\smallskip}}
